\section{Hamiltonian Large Language Models (HLLM)}


\subsection{Core Concept}

Traditional fine-tuning modifies model parameters $\theta$ via gradient descent:
\[
\theta \leftarrow \theta - \alpha \nabla_\theta \mathcal{L}(\theta, \mathcal{D})
\]
This is expensive (requires backpropagation through billions of parameters) and opaque (weight changes are inscrutable).

\textbf{HLLM} \cite{zoo2025hllm} proposes an alternative: keep $\theta$ frozen and optimize the \textit{context} $\Psi$ instead. The system obeys a \textbf{Hamiltonian invariant}:
\[
\Psi \cdot \Theta = \kappa
\]
where:
\begin{itemize}
\item $\Psi$ = Policy mass (semantic context / experiences)
\item $\Theta$ = Inference cost (model entropy / uncertainty)
\item $\kappa$ = Conserved constant (system equilibrium)
\end{itemize}

As context expands ($\Psi \uparrow$), model uncertainty decreases ($\Theta \downarrow$) while preserving total system cost $\kappa$.

\subsection{Training-Free GRPO}

\textbf{Group Relative Policy Optimization (GRPO)} \cite{shao2024deepseekmath} computes advantages relative to group means instead of baselines:
\[
A_i = \frac{R_i - \mu_G}{\sigma_G}
\]
where $R_i$ is the reward for output $i$, $\mu_G$ is the group mean, and $\sigma_G$ is group standard deviation.

\textbf{Training-Free GRPO} \cite{tencent2025grpo} replaces numerical advantages with \textit{semantic advantages}—natural language insights extracted via LLM introspection. The algorithm proceeds in three stages:

\subsubsection{Stage 1: Trajectory Summarization}

For each generated output $o_i$, an LLM summarizes:
\begin{itemize}
\item What actions were taken
\item Which experiences were used
\item Where errors or detours occurred
\end{itemize}

\begin{algorithm}[H]
\caption{Trajectory Summarization}
\begin{algorithmic}[1]
\STATE \textbf{Input:} Trajectory $\tau_i$, correctness label $c_i$, ground truth $y^*$
\STATE \textbf{Output:} Natural language summary $s_i$
\STATE $s_i \leftarrow \text{LLM.summarize}(\tau_i, c_i, y^*)$
\RETURN $s_i$
\end{algorithmic}
\end{algorithm}

\subsubsection{Stage 2: Group Advantage Extraction}

Given $G$ trajectory summaries for a query, the LLM identifies patterns distinguishing successes from failures. It outputs up to 3 operations per group: \texttt{add}, \texttt{modify}, or \texttt{delete} experiences.

\begin{algorithm}[H]
\caption{Group Advantage Extraction}
\begin{algorithmic}[1]
\STATE \textbf{Input:} Summaries $\{s_1, \ldots, s_G\}$, current experiences $E$, ground truth $y^*$
\STATE \textbf{Output:} Experience operations $\mathcal{O}_G$
\STATE $\mathcal{O}_G \leftarrow \text{LLM.extract\_advantages}(\{s_i\}, E, y^*)$
\STATE Filter to max 3 operations: $|\mathcal{O}_G| \leq 3$
\RETURN $\mathcal{O}_G$
\end{algorithmic}
\end{algorithm}

\subsubsection{Stage 3: Batch Consolidation}

All group operations from a batch are consolidated to avoid duplication and ensure experiences are $\leq 32$ words.

\begin{algorithm}[H]
\caption{Batch Consolidation}
\begin{algorithmic}[1]
\STATE \textbf{Input:} All group operations $\{\mathcal{O}_1, \ldots, \mathcal{O}_B\}$, current experiences $E$
\STATE \textbf{Output:} Final consolidated operations $\mathcal{O}_{\text{final}}$
\STATE $\mathcal{O}_{\text{final}} \leftarrow \text{LLM.consolidate}(\{\mathcal{O}_j\}, E)$
\STATE Apply operations: $E \leftarrow E \oplus \mathcal{O}_{\text{final}}$
\RETURN Updated experience library $E$
\end{algorithmic}
\end{algorithm}

\subsection{Performance}

On the AIME mathematical reasoning benchmark \cite{hendrycks2021math}, Training-Free GRPO achieves:
\begin{itemize}
\item \textbf{AIME24}: 82.7\% (vs. 80.0\% vanilla GRPO, +2.7\%)
\item \textbf{AIME25}: 73.3\% (vs. 67.9\% vanilla GRPO, +5.4\%)
\item \textbf{Training Cost}: \$18 (vs. \$10,000+ for fine-tuning)
\item \textbf{Training Samples}: 100 (vs. 10,000+ for fine-tuning)
\end{itemize}

Crucially, base model weights remain \textit{frozen}—all improvement comes from context optimization.

